\section{Implementation}

\subsection{Project Structure}
AgriApp is organized as a mono-repository where backend, web UI, Android control app, firmware helpers, and operational scripts evolve together. This structure is intentional: the project is deployed as a single edge software stack, and versioning the full stack in one place reduces integration drift.

\begin{verbatim}
agri-app/
|-- backend/                  # Python backend helpers (capture/sensors)
|-- docs/
|   `-- software-doc-latex/   # Technical documentation (this report)
|-- firmware/
|   `-- esp32-cam/            # ESP32-CAM firmware + utilities
|-- mobile/
|   `-- android/              # Android tablet app (Kotlin + Compose)
|-- old/                      # Deprecated scripts kept for reference
|-- scripts/                  # Runtime/deploy helper scripts
|   |-- capture_rpi.sh
|   |-- capture_esp.py
|   |-- capture_pair.sh
|   |-- esp_capture.sh
|   |-- rpi_network_mode.sh
|   |-- run_server.sh
|   `-- install_systemd_user_services.sh
|-- static/
|   |-- css/                  # Web styles
|   |-- js/                   # Web logic
|   `-- uploads/              # images/ labels/ jsons/ metadata/ thumbs/
|-- systemd/                  # user service templates
|-- templates/                # Flask HTML templates
|-- test/                     # Hardware validation scripts (GPS/BME280)
|-- client.py                 # Raspberry Pi camera capture client
|-- server.py                 # Flask backend + API orchestration
|-- deploy_rpi.sh             # Sync/mount/reload/android/esp helper
`-- README.md
\end{verbatim}

\subsection{Data Model and Filesystem Invariants}
The runtime data model is intentionally file-oriented and must preserve strict naming coherence across sibling folders. For a given image stem (for example \texttt{rpi\_20260408-152834}), expected artifacts are:
\begin{itemize}
  \item \texttt{images/<stem>.jpg} (or compatible image extension),
  \item \texttt{metadata/<stem>.json} (sensor/context sidecar),
  \item optional \texttt{jsons/<stem>.json} (annotation JSON),
  \item optional \texttt{labels/<stem>.txt} (YOLO export),
  \item optional \texttt{masks/<stem>.png} (SegGPT/segmentation result mask),
  \item optional \texttt{thumbs/<stem>.jpg} (cache artifact).
\end{itemize}

Two invariants are operationally important:
\begin{enumerate}
  \item delete operations are lifecycle-complete: deleting one image must remove related sidecars/caches to avoid orphan artifacts;
  \item metadata persistence is best-effort but deterministic in key set: missing sensors produce \texttt{null} values, not schema drift.
\end{enumerate}

For segmentation-assisted Studio workflows, \texttt{masks/\_last\_prompt.png} is a reusable global prompt artifact. It is not image-specific and can be overwritten or removed without affecting captured binaries.

\subsection{File Lifecycle Matrix (Create/Update/Delete)}
\begin{longtable}{p{0.26\linewidth} p{0.22\linewidth} p{0.44\linewidth}}
\toprule
Operation & Target & Effects on Artifacts \\
\midrule
RPi one-shot capture & \texttt{images/<stem>.jpg} & Creates image; creates/updates \texttt{metadata/<stem>.json}; no labels by default. \\
ESP one-shot capture (proxy) & \texttt{images/esp\_<stem>.jpg} & Creates image; metadata sidecar best-effort; normalization can backfill missing sidecar. \\
Label save & \texttt{jsons/<stem>.json}, \texttt{labels/<stem>.txt} & Rewrites both JSON and YOLO TXT from current in-memory label set. \\
SegGPT oneshot (Studio) & \texttt{masks/<stem>.png}, label files & Writes segmentation mask and updates annotation files through same label save path. \\
Last mask update & \texttt{masks/\_last\_prompt.png} & Overwrites reusable prompt-mask artifact. \\
Delete single image & stem-linked files & Removes image and associated \texttt{jsons/labels/metadata/thumbs/masks} artifacts for the same stem. \\
Delete all images & upload subtree data & Clears all capture-related artifacts in managed upload folders. \\
Dataset import (Studio) & \texttt{images/labels/jsons/metadata} & Adds files; optional overwrite policy for collisions. \\
\bottomrule
\end{longtable}

\subsection{Backend Implementation}
The backend is implemented in Python/Flask and runs on Raspberry Pi as the authoritative runtime node. The implementation is API-centric: business logic for capture, gallery lifecycle, annotations, network mode switching, and system operations is concentrated server-side, while clients (web/tablet) remain thin.

In concrete terms, \texttt{server.py} currently implements five operational domains. The first is acquisition orchestration: one-shot capture from Raspberry Pi camera through \texttt{scripts/capture\_rpi.sh}, and ESP proxy capture through \texttt{scripts/capture\_esp.py}. The second is gallery lifecycle management: paginated listing, per-image status, thumbnail generation/caching, single delete, and delete-all operations. The third is labeling persistence and conversion to dual artifacts (JSON + YOLO TXT). The fourth is device control and connectivity (restart/reboot/poweroff, Wi-Fi/AP mode switching). The fifth is unattended acquisition control through start/stop/status APIs for capture loops. The loop implementation is split by source, with dedicated RPi and ESP loop states and endpoints.

An implementation refinement introduced during field testing concerns ESP unattended capture: the ESP loop now triggers the internal ESP proxy endpoint rather than writing images directly through a standalone script invocation. This guarantees that each ESP image follows the same persistence path as one-shot captures (image save + metadata enrichment + sidecar generation), reducing divergence between manual and unattended workflows.

A notable implementation choice is best-effort metadata enrichment. After each successful capture, the backend attempts to enrich the image with GPS and BME280 values and writes a sidecar JSON. Capture completion is never blocked by missing sensors: if sensors are unavailable, the image remains valid and metadata fields are populated as \texttt{null}. For ESP artifacts, metadata sidecars are additionally normalized at gallery-read time, so missing sidecars from legacy or transitional runs are automatically repaired.

Recent hardening focused on GPS serial access under concurrent API traffic. The GPS reader now uses a process-level lock and short TTL cache, so repeated calls (for example health checks plus capture paths) do not repeatedly open/read the UART in parallel. This significantly reduces contention symptoms on \texttt{/dev/serial0} and stabilizes metadata enrichment during normal operations.

\subsection{Web Frontend Implementation}
The web frontend is implemented with Flask templates plus static CSS/JavaScript modules. Its architecture separates gallery concerns from labeling concerns, reducing coupling and allowing iterative UI work without destabilizing annotation behavior.

Recent refactoring introduced a shared template shell (\texttt{base.html}) and reusable partials for navbar/footer. This removed duplicated header/footer markup across pages (\texttt{index}, \texttt{gallery}, \texttt{capture}, \texttt{system}, \texttt{health}, \texttt{log}) and reduced visual drift between views. A practical consequence is faster UI iteration with lower regression risk when navigation or branding elements change.

To improve field robustness, frontend dependencies were moved to locally served static assets (Bootstrap bundle and Socket.IO client), and external web-font dependencies were removed in favor of system-font stacks. This allows the web UI to render correctly even in isolated AP-only deployments where Internet access is unavailable.

The gallery view is driven by server-side pagination and lightweight thumbnails. This is critical for large datasets: the list/grid view uses cached reduced images, while full-resolution images are loaded only in detail workflows. Mutation operations (save labels, delete image, delete-all) are followed by explicit API synchronization instead of optimistic local state assumptions.

The labeler view is implemented as an interaction-focused editor for bounding boxes and class assignment. Persistence is dual-output by design: JSON preserves full UI state, while TXT keeps compatibility with YOLO pipelines.

The \texttt{System} page and \texttt{Gallery} destructive flows now use custom confirmation modals instead of browser default popups, aligning interaction style with the rest of the UI and making high-impact operations (power actions, delete-all) more explicit.
In the current field baseline, the web UI is intentionally exposed in RPi-only control mode: ESP controls and ESP status blocks are hidden from \texttt{/capture}, \texttt{/system}, and \texttt{/health} to keep operations simpler and less error-prone. ESP proxy endpoints are still present backend-side for optional diagnostics and controlled experiments.

\subsection{Android App Implementation}
The Android application is written in Kotlin with Jetpack Compose and Retrofit. The UI is state-driven through a central view-model that coordinates discovery, control actions, gallery pagination, and logs.

The app currently exposes two major control planes. The first is connectivity and device operation: discovery, Wi-Fi/AP switching, capture start/stop, one-shot actions for \texttt{RPi}/\texttt{ESP}/\texttt{Both}, and system actions (restart/reboot/poweroff). Recent iterations moved capture and system controls from popup-driven flows to inline controls, reducing interaction friction on phone portrait layouts. Auto-capture start supports source selection (\texttt{RPi}, \texttt{ESP}, or \texttt{Both}) with interval presets (including 30s and minute-scale options), and UI guardrails keep the flow unambiguous: start is disabled while loops are active, while stop is enabled only when relevant loops are running. In addition, auto-loop status synchronization at app bootstrap now aligns UI source/interval chips with backend loop state, avoiding stale selection markers after app restart.

The second control plane is gallery consumption, implemented with paginated grid views and thumbnail-first loading. Each card shows compact source-aware timestamp formatting and contextual metadata summaries. Image detail is handled in a dedicated full-screen activity optimized for portrait devices: zoom/pan gestures, explicit \texttt{Prev}/\texttt{Next} navigation buttons, destructive delete with confirmation, and multiline metadata rendering (GPS, temperature, humidity, pressure) together with image resolution and file size. This separation keeps list interaction lightweight while preserving rich inspection workflows in detail view.

A practical implementation detail is cache-first discovery. Last known hosts are persisted locally and probed before subnet scanning. The current discovery policy is intentionally deterministic: RPi probing follows the current client subnet, while ESP probing is constrained to \texttt{192.168.4.0/24} (rescue/AP segment), avoiding ambiguous hybrid scans. During network mode switches (Wi-Fi $\leftrightarrow$ AP), the app triggers an automatic discovery step after a short delay to recover quickly from expected subnet transitions. Another relevant detail is event-driven in-app logging: connection transitions, discovery outcomes, auto-loop lifecycle events, and shot outcomes are surfaced in a bottom log panel using concise operational messages.

The \texttt{Connections} and \texttt{System} tabs were also restructured for operational clarity. \texttt{Connections} now emphasizes quick reachability context (host chips, AP/Wi-Fi ON/OFF badges, and per-device ``last seen'' indicators). \texttt{System} is segmented into dedicated cards (\texttt{Storage}, \texttt{WiFi}, \texttt{Server}, \texttt{Power}) so high-impact actions are visually separated from network-selection tasks. The Wi-Fi card supports manual ``connect now'' against preconfigured NetworkManager profiles, enabling deterministic network switching from the mobile client.

\subsection{Storage Implementation}
Persistence is filesystem-based and explicit:
\begin{itemize}
  \item \texttt{static/uploads/images}: original captured images,
  \item \texttt{static/uploads/labels}: YOLO TXT labels,
  \item \texttt{static/uploads/jsons}: full annotation state,
  \item \texttt{static/uploads/metadata}: per-image environmental/geospatial metadata,
  \item \texttt{static/uploads/thumbs}: cached generated thumbnails.
\end{itemize}

The implementation enforces lifecycle consistency: delete endpoints remove associated label/json/metadata artifacts and thumbnail cache entries. Capture naming also encodes source semantics (e.g., \texttt{rpi\_...} and \texttt{esp\_...}) to simplify mixed-source campaigns.

\subsection{Operational Script Inventory}
The following scripts are currently considered operationally relevant for reproducible deployments and evaluations:
\begin{longtable}{p{0.28\linewidth} p{0.67\linewidth}}
\toprule
Script & Role \\
\midrule
\texttt{deploy\_rpi.sh} & Unified operator entrypoint: sync, mount, remote commands, server reload, ESP build/flash, Android build/install/run. \\
\texttt{scripts/run\_server.sh} & Backend launcher with environment bootstrapping (\texttt{.env.systemd}, pyenv/venv). \\
\texttt{scripts/capture\_rpi.sh} & Raspberry Pi capture wrapper (oneshot and interval loop modes). \\
\texttt{scripts/capture\_esp.py} & ESP helper (HTTP and serial capture modes). \\
\texttt{scripts/rpi\_network\_mode.sh} & NetworkManager integration for Wi-Fi/AP mode switching. \\
\texttt{scripts/seggpt\_compare\_benchmark.py} & Cross-node SegGPT benchmark (Studio vs RPi), CSV/JSON artifact export. \\
\texttt{scripts/install\_systemd\_user\_services.sh} & RPi user-service installation and enablement. \\
\texttt{scripts/install\_local\_user\_service.sh} & Studio local service installation (\texttt{agriapp-local.service}). \\
\texttt{test/test-gps.py} & UART-level GPS diagnostic utility. \\
\texttt{test/test-bme.py} & I2C-level BME280 diagnostic utility. \\
\bottomrule
\end{longtable}

\subsection{Implementation Trade-offs}
The software intentionally favors operational clarity over abstraction depth. A single backend process on Raspberry Pi simplifies deployment and observability, at the cost of centralizing responsibility on the edge node. Filesystem persistence eases backup and manual inspection, at the cost of weaker ad-hoc indexing compared to database-backed designs. These trade-offs are acceptable for the current project phase, where field robustness and rapid iteration are the primary objectives.
